В первой главе рассматриваются математические основы интерполяции на разреженных сетках, лежащие в основе разработанного вычислительного ядра. Основное внимание уделяется иерархическому представлению аппроксимируемой функции, построению базисных функций и критериям адаптивного уточнения. Отдельно обсуждается применение полученного аппарата к аппроксимации операторов, возникающих при численном решении дифференциальных уравнений.

\subsection{Интерполяция многомерных функций и проблема роста вычислительной сложности}

Пусть задано отображение
\begin{equation}
    F \colon \Omega \subset \mathbb{R}^d \to \mathbb{R}^m,
\end{equation}
которое требуется многократно вычислять на области $\Omega = [a_1, b_1] \times \dots \times [a_d, b_d]$. Если каждое вычисление $F(x)$ связано с численным интегрированием, решением вспомогательной задачи или интерпретацией сложной формулы, возникает задача построения приближённого интерполянта
\begin{equation}
    \widetilde{F}(x) \approx F(x), \quad x \in \Omega,
\end{equation}
который обеспечивал бы достаточно малую ошибку при существенно меньшей стоимости вычисления.

Прямое построение полной тензорной сетки в $d$-мерном случае требует экспоненциального роста числа узлов. Если в каждом измерении используется $n$ узлов, то общее число точек равно $n^d$. Такое поведение известно как «проклятие размерности». Даже для умеренных значений $n$ и $d$ этот подход становится неэффективным. Поэтому вместо полной сетки применяются разреженные конструкции, в которых сохраняются только наиболее информативные комбинации одномерных уровней детализации.

\subsection{Нормировка области и постановка задачи интерполяции}

В разработанной системе вычисления выполняются на единичном гиперкубе. Для этого исходная точка $x = (x_1, \dots, x_d)$ переводится из реальной области $\Omega$ в нормированные координаты $u = (u_1, \dots, u_d)$ по формуле
\begin{equation}
    u_i = \frac{x_i - a_i}{b_i - a_i}, \quad i = 1, \dots, d.
\end{equation}
Обратное преобразование имеет вид
\begin{equation}
    x_i = a_i + u_i (b_i - a_i).
\end{equation}

Такое нормирование позволяет единообразно задавать базисные функции, не зависящие от конкретного диапазона аргументов. В результате задача аппроксимации переносится на область $[0,1]^d$, а геометрия сетки определяется только уровнями и индексами узлов.

\subsection{Одномерные иерархические базисные функции}

Построение разреженной сетки основано на иерархическом базисе. В одномерном случае каждому узлу сопоставляется уровень $\ell$ и индекс $i$. Для граничного уровня $\ell = 0$ используются две базисные функции, отвечающие концам отрезка:
\begin{equation}
    \varphi_{0,0}(x) = 1 - x, \qquad \varphi_{0,1}(x) = x.
\end{equation}
Эти функции обеспечивают корректную интерполяцию на концах области и задают начальный слой разложения.

Для уровней $\ell > 0$ положение центра узла определяется как
\begin{equation}
    c_{\ell,i} = \frac{i}{2^{\ell}},
\end{equation}
а радиус влияния равен
\begin{equation}
    h_{\ell} = \frac{1}{2^{\ell}}.
\end{equation}
Если аргумент находится вне области влияния узла, вклад базисной функции равен нулю. Иначе базис зависит от выбранного типа аппроксимации.

Для линейного базиса используется выражение
\begin{equation}
    \varphi_{\ell,i}^{(lin)}(x) = 1 - \frac{|x - c_{\ell,i}|}{h_{\ell}}, \quad |x-c_{\ell,i}| < h_{\ell}.
\end{equation}

Для квадратичного базиса, который является основным в данной работе, применяется функция
\begin{equation}
    \varphi_{\ell,i}^{(quad)}(x) = 1 - \left(\frac{|x - c_{\ell,i}|}{h_{\ell}}\right)^2, \quad |x-c_{\ell,i}| < h_{\ell}.
\end{equation}

Квадратичный базис обеспечивает более гладкое изменение интерполянта внутри области влияния узла и позволяет лучше аппроксимировать гладкие отображения при той же структуре уровней. Это особенно важно при визуализации динамических систем, где требуется избегать изломов и искусственных артефактов.

\subsection{Многомерный базис и тензорное произведение}

Для точки с многомерными координатами $u = (u_1, \dots, u_d)$ вклад узла с уровнями $\ell = (\ell_1, \dots, \ell_d)$ и индексами $i = (i_1, \dots, i_d)$ задаётся произведением одномерных функций:
\begin{equation}
    \Phi_{\ell,i}(u) = \prod_{k=1}^{d} \varphi_{\ell_k,i_k}(u_k).
\end{equation}

Такое определение соответствует тензорному произведению одномерных базисов. Если хотя бы в одном измерении аргумент выходит из области влияния соответствующего узла, то весь многомерный базис обращается в нуль. Благодаря этому вычисление интерполянта носит локальный характер: в ответе участвуют только те узлы, чьи области поддержки пересекают рассматриваемую точку.

\subsection{Иерархическое представление и коэффициенты избытка}

Интерполянт на разреженной сетке строится в виде суммы иерархических вкладов:
\begin{equation}
    \widetilde{F}(u) = \sum_{(\ell,i) \in \mathcal{G}} \alpha_{\ell,i} \, \Phi_{\ell,i}(u),
\end{equation}
где $\mathcal{G}$ --- множество узлов сетки, а $\alpha_{\ell,i} \in \mathbb{R}^m$ --- иерархические коэффициенты, или коэффициенты избытка.

Смысл коэффициента $\alpha_{\ell,i}$ состоит в том, что он измеряет новую информацию, вносимую узлом текущего уровня по сравнению с уже построенным приближением. Если обозначить через $u_{\ell,i}$ координаты центра узла, то коэффициент вычисляется как разность между точным значением и уже имеющейся интерполяцией:
\begin{equation}
    \alpha_{\ell,i} = F(u_{\ell,i}) - \widetilde{F}_{prev}(u_{\ell,i}).
\end{equation}

Если норма коэффициента мала, это означает, что в данной точке текущее приближение уже хорошо описывает поведение функции и дальнейшее уточнение может быть нецелесообразно. Наоборот, большие значения $\|\alpha_{\ell,i}\|$ указывают на необходимость детализации сетки.

\subsection{Структура разреженной сетки и входные точки}

В разреженной сетке каждый узел задаётся набором уровней и индексов по всем измерениям. Уровень определяет масштаб локализации, а индекс --- положение центра. В отличие от полной тензорной сетки, в разреженной конструкции рассматриваются только те комбинации уровней, которые удовлетворяют ограничению на суммарную сложность. В реализованном алгоритме начальные входные точки строятся как комбинации активных измерений, после чего для каждой такой конфигурации выполняется локальное адаптивное уточнение.

На практике это означает, что сетка состоит из семейства входных узлов, соответствующих различному числу «активированных» координатных направлений. Такое представление удобно как с математической точки зрения, так и с точки зрения реализации рекурсивного обхода структуры при вычислении интерполянта.

\subsection{Адаптивное уточнение и локальный критерий продолжения}

Одним из ключевых преимуществ разреженных сеток является возможность адаптивного уточнения. В разработанном алгоритме решение о создании дочерних узлов принимается на основе двух факторов:
\begin{enumerate}
    \item нормы иерархического коэффициента текущего узла;
    \item поведения ошибки в специальных контрольных точках.
\end{enumerate}

Если для текущего узла выполняется условие
\begin{equation}
    \|\alpha_{\ell,i}\| \leq \varepsilon,
\end{equation}
то узел может считаться достаточно точным. Однако этого недостаточно для надёжной аппроксимации функций, имеющих нули или симметричные участки. Поэтому дополнительно рассматриваются контрольные точки, или anchor points. Если внутри области влияния узла имеется контрольная точка $u^{(a)}$, для которой выполняется
\begin{equation}
    \left\|F\left(u^{(a)}\right) - \widetilde{F}\left(u^{(a)}\right)\right\| \geq \varepsilon,
\end{equation}
уточнение продолжается даже в случае малого локального коэффициента. Такой механизм особенно важен для функций типа $\sin(x)$, где значения на концах и в центре могут не выявлять внутреннюю структуру.

В результате адаптивный алгоритм концентрирует узлы в областях с высокой кривизной, быстрым изменением или плохо учтённой динамикой, а на гладких участках оставляет более грубое разбиение.

\subsection{Порождение дочерних узлов}

Пусть текущий узел имеет уровень $\ell_k > 0$ в $k$-м измерении. Тогда при уточнении по этой координате строятся два дочерних узла с уровнем
\begin{equation}
    \ell_k' = \ell_k + 1,
\end{equation}
и индексами
\begin{equation}
    i_k^{left} = 2 i_k - 1, \qquad i_k^{right} = 2 i_k + 1.
\end{equation}
Все остальные уровни и индексы сохраняются. Такая формула отражает стандартное иерархическое ветвление: каждый внутренний узел порождает левое и правое уточнение относительно своего центра.

При этом фактическое создание потомков в реализации зависит от направления ошибки. Если контрольная точка расположена левее центра узла, имеет смысл прежде всего уточнять левую ветвь, и наоборот. Это снижает число лишних вычислений и уменьшает размер итоговой сетки.

\subsection{Рекурсивное вычисление интерполянта}

После построения сетки значение интерполянта в точке $u$ вычисляется как сумма вкладов входных узлов и их потомков. Процесс удобно описывается рекурсивно. Для узла $n$ с коэффициентом $\alpha_n$ и базисной функцией $\Phi_n$ локальный вклад равен
\begin{equation}
    V_n(u) = \alpha_n \Phi_n(u).
\end{equation}
Если узел имеет потомков, к этому значению добавляются их рекурсивные вклады:
\begin{equation}
    S_n(u) = V_n(u) + \sum_{c \in Children(n)} S_c(u).
\end{equation}

Благодаря локальности базисных функций рекурсия быстро обрывается: если $\Phi_n(u)=0$, вклад всей соответствующей ветви становится нулевым. Это делает вычисление интерполянта существенно дешевле, чем прямой пересчёт исходной функции.

\subsection{Аппроксимация отображений и композиция операторов}

Разработанный метод применим не только к обычным функциям вида $F(x)$, но и к операторам перехода. Пусть задано отображение
\begin{equation}
    G \colon \Omega \to \Omega,
\end{equation}
описывающее один шаг эволюции системы. Тогда можно построить интерполянт $\widetilde{G}$ и затем использовать его как аргумент для следующего оператора. Если требуется аппроксимировать композицию двух отображений $F \circ G$, то естественно использовать формулу
\begin{equation}
    (F \circ G)(x) = F(G(x)).
\end{equation}

В реализации это соответствует построению первой сетки для $G$, после чего вычисление новой сетки выполняется уже для отображения, принимающего на вход результат интерполяции $\widetilde{G}(x)$. Такой подход позволяет моделировать последовательную эволюцию системы по шагам без повторного хранения всей траектории.

\subsection{Связь с задачами решения дифференциальных уравнений}

Пусть динамическая система задаётся обыкновенным дифференциальным уравнением
\begin{equation}
    \frac{dz}{dt} = f(z,t), \qquad z(t_0) = z_0.
\end{equation}
Тогда оператор перехода на временном шаге $\Delta t$ можно записать как отображение
\begin{equation}
    \Psi_{\Delta t}(z_0) = z(t_0 + \Delta t).
\end{equation}

Если для вычисления $\Psi_{\Delta t}$ используется численный интегратор, то сам оператор может быть дорогим. В таком случае имеет смысл построить его интерполяцию на разреженной сетке. После этого задача многократного интегрирования заменяется многократным вычислением интерполянта. Для нескольких последовательных временных шагов применяется композиция операторов:
\begin{equation}
    \Psi_{k \Delta t} = \underbrace{\Psi_{\Delta t} \circ \dots \circ \Psi_{\Delta t}}_{k \text{ раз}}.
\end{equation}

Именно эта идея лежит в основе поддержки дифференциальных уравнений в проекте: в качестве аппроксимируемой функции рассматривается либо непосредственное численное решение на заданном промежутке времени, либо одношаговый оператор, который затем композиционно повторяется.

\subsection{Преимущества квадратичного базиса в рассматриваемой задаче}

Хотя линейный базис проще и традиционно широко используется, в данной работе основной акцент сделан на квадратичном базисе. Его применение оправдано следующими причинами:
\begin{enumerate}
    \item квадратичная форма лучше описывает гладкие участки отображения;
    \item при одинаковой глубине сетки достигается меньшая ошибка на плавных функциях;
    \item для визуализации векторных полей и фазовых портретов уменьшается число заметных изломов между соседними областями влияния;
    \item квадратичный базис естественно согласуется с идеей адаптивного уточнения, позволяя уточнять только действительно сложные зоны.
\end{enumerate}

При этом архитектура реализации сохраняет возможность использования и линейного базиса, что удобно для сравнения и бенчмаркинга.

\subsection*{Выводы по главе}
\addcontentsline{toc}{subsection}{Выводы по главе}

В главе рассмотрены математические основы метода, используемого в дипломной работе. Показано, что переход к разреженным сеткам позволяет уменьшить число узлов по сравнению с полной тензорной дискретизацией, а иерархическое представление на базе коэффициентов избытка обеспечивает естественный механизм адаптивного уточнения. Обосновано использование квадратичного базиса как более точного инструмента для аппроксимации гладких отображений. Также показано, что аппарат разреженных сеток применим не только к обычным функциям, но и к операторам перехода, возникающим при численном решении дифференциальных уравнений. Эти положения образуют теоретическую базу для архитектуры и алгоритмов, реализованных в программном комплексе.